Hospitals accumulate vast electronic health record (EHR) data that remains
largely unprocessed and inaccessible to the non-technical staff who could act
on it. We present \textit{HospitAI}, a healthcare analytics platform that turns
raw EHR data into clinician-readable insight through a single, layered, mostly
unsupervised pipeline. Raw clinical events from more than ten heterogeneous
tables are ingested and reduced to one feature vector per patient mixing
numerical and categorical attributes. Patients are then grouped with the
K-Prototypes algorithm, which handles mixed-type data within one distance
function. On top of the resulting groups we run a \emph{dual-scope}
sequential pattern miner: PrefixSpan is applied both globally and
\emph{independently within each cluster}, so that care trajectories
characteristic of a patient group can be separated from population-wide ones.
A large language model (LLM) layer sits across every analytical step: it
converts statistical cluster profiles into plain-language summaries, interprets
mined clinical sequences, and exposes a natural-language text-to-SQL interface
governed by a validate-then-execute state machine. Using a synthetic cohort of
\NPATIENTS{} patients generated with Synthea, we obtain interpretable patient
clusters and cluster-specific clinical trajectories, and we show how the LLM
layer makes these results consumable by users without data-analysis training.
The contribution is not any single algorithm but their integration into an
end-to-end, interpretation-first analytics pipeline.
